\documentclass[handout]{beamer}
%\documentclass[presentation]{beamer}

% -----------------------------------------------------------------------------
% Tema do beamer
\usecolortheme{ImeccRed}

% -----------------------------------------------------------------------------
% Pacotes básicos 
\usepackage[utf8]{inputenc}
\usepackage[portuges]{babel}
\usepackage{booktabs}
\usepackage{caption}
\usepackage{subcaption}
\usepackage{graphicx}
\usepackage{amsmath}
\usepackage{amsfonts}
\usepackage{amssymb}
\usepackage{epstopdf}
\usepackage{csquotes}
\usepackage{url}
\usepackage[natbib=true,style=authoryear, backend=biber, useprefix=true]{biblatex}
\addbibresource{bibliografia.bib}

% -----------------------------------------------------------------------------
% Pacotes (Outros) - acrescente os pacotes que você precisa abaixo!


% -----------------------------------------------------------------------------
% Logo do Imecc + Unicamp estilizada.
\institute{\includegraphics[scale=.35]{imecc}}

% -----------------------------------------------------------------------------
%Informacoes para incluir na primeira página:
\title{Título}

\subtitle{Subtítulo}

\author{Autor(a)}

\date{\today}

% -----------------------------------------------------------------------------
\begin{document}
	
\frame{\titlepage}
	
    \begin{frame}
		\frametitle{Tabela e bloco de texto}
		\begin{table}[h]
            \centering
            \begin{tabular}{ccc}
                \hline
                Coluna 1  & Coluna 2 & Coluna 3 \\ 
                \hline
                texto   & texto   & texto  \\
                texto   & texto   & texto  \\
                texto   & texto   & texto  \\
                \hline
            \end{tabular}\label{tabela01}
        \end{table}
		
		\begin{block}{Título do Bloco}
			Texto Texto Texto Texto Texto
			
			Testando a referência, (\cite{tantau2004user}).
		\end{block}
    \end{frame}
	
	
    \begin{frame}
		\frametitle{Figura lado a lado}
		\begin{columns}[b]
			\column{0.0125\textwidth}
			\column{0.4875\textwidth}
			\centering
			\begin{figure}
				\includegraphics[width=0.8\textwidth]{figA.jpeg} \
				\caption{Figura A. Fonte: \cite{google}.}
			\end{figure}
			\column{0.4875\textwidth}
			\centering
			\begin{figure}
				\includegraphics[width=0.8\textwidth]{figB.jpeg} \
				% remove the 'draft' keyword, when replacing with final figure!
				\caption{Figura B. Fonte: \cite{google}.}
			\end{figure}
			\column{0.0125\textwidth}
		\end{columns}
	\end{frame}
	
	
	\begin{frame}
		\frametitle{Figura e texto lado a lado}
		\begin{columns}[]
			\column{0.0125\textwidth}
			\column{0.4875\textwidth}
			\centering
			\begin{figure}
				\includegraphics[width=0.8\textwidth]{figC.jpeg} \
				\caption{Figura C. Fonte: \cite{google}.}
			\end{figure}
			\column{0.4875\textwidth}
			Texto Texto Texto Texto Texto 
			\begin{itemize}
				\item Item 1
				\item Item 2
				\item Item 3
				\item Item 4
			\end{itemize}			
			\column{0.0125\textwidth}
		\end{columns}
	\end{frame}
	
	
	\begin{frame}
		\frametitle{Uma figura}
		\bigskip
		\begin{figure}
			\includegraphics[width=0.6\textwidth, height=0.4\textwidth]{figD}
			\caption{Figura D. Fonte: \cite{google}.}
		\end{figure}
	\end{frame}

% -----------------------------------------------------------------------------
% Referências - você pode adicionar suas referências no arquivo bibliografia.bib

\begin{frame}[t, allowframebreaks]{}
	\frametitle{Referências}
    \printbibliography
\end{frame}
	
	
\end{document}